% =====================================================================
\section{Conclusions and future work}\label{sec:conclusions}
% =====================================================================

We have proposed Rate-Independent Epigenetics (RIE), a framework that
models epigenetic change as rate-independent dissipative evolution
driven by a single triple $(\Q,E,\Psi)$: an energy landscape of
chromatin configurations, a micro-environmental loading, and a
$1$-homogeneous dissipation potential encoding the resistance to
remodelling. Once the triple is fixed and an energetic principle is
postulated, the governing equations---the energetic formulation and the
equivalent Biot inclusion---follow systematically, and their solutions
satisfy, a priori, estimates that coincide with the first and second
laws of thermodynamics. The $1$-homogeneity of the dissipation
potential is the structural source of rate independence: it yields the
saturation identity, the non-negativity and minimality of the
dissipation, and the purely geometric characterisation of the
hysteretic dissipation over closed cycles. Threshold activation,
memory, and residual state upon unloading---the defining phenomenology
of epigenetic marks---are thus not modelled by ad-hoc equations but
emerge from a thermodynamically consistent variational structure.

On the analytical side we established the well-posedness of RIE models,
making explicit the compatibility argument behind existence and adding
the definiteness hypothesis needed for the balanced-viscosity
construction. Energetic solutions exist under coercivity and
continuity; they are unique under convexity and, for the scalar
bistable landscapes of epigenetic interest, the balanced-viscosity
selection restores uniqueness without convexity. On the numerical side
we analysed an incremental variational integrator: it is
unconditionally stable in the energetic sense, converges to energetic
solutions, and is first-order globally consistent in the energy
balance, reducing in the scalar case to a two-sided return map. The
linear play operator---the minimal RIE model of epigenetic
memory---admits a closed-form evolution exhibiting both reversible and
irreversible responses, and the integrator reproduced it exactly at the
nodes, confining the discretisation error to the energy balance.

Several directions remain open. The balanced-viscosity theory is
developed here for scalar states; its extension to vectorial
epigenetic variables, where the transition path at a switch is
genuinely path-dependent, is an important next step, as is a rigorous
convergence rate for the state in the presence of switches. The
local truncation estimate invites adaptive time-stepping, and
non-symmetric dissipation potentials would model the
direction-dependent resistance of methylation versus demethylation.

Most importantly, the present framework is the abstract foundation for
a quantitative epigenetic model. In a companion paper we specialise RIE
to a scalar \emph{bistable} landscape---a double-well free energy with a
saturating micro-environmental coupling---in which the healthy and
damaged phenotypes are the two wells. There the convexity of the
present benchmark is lost, catastrophic phenotype switches (of
epithelial--mesenchymal type) appear as energetic jumps selected by the
balanced-viscosity criterion, and the analytical thresholds, residual
``epigenetic scars'' and bifurcation structure can be calibrated
against experimental data on hypoxia-driven epigenetic remodelling.
